\chapter{Этаж~2. Составной узел: кварки и квантовые числа}
\label{ch:floor2}

\section{Радиус и объект}
\label{sec:floor2-scale}

Этаж~2 на оси длины (\cref{prop:ladder-scale-steps}):
\[
  N\sim 10^{15}\cdot\hl,\qquad \log_{10}N\approx 15.
\]
Это уже \emph{составной узел}: несколько фокусов планконов в одной $\varepsilon$-окрестности,
связанных топологией FCC и запретом Паули (A16), а не один vortex на одном~$\hv$.

На~$T$ здесь --- схема конфайнмента и кварковых степеней свободы;
на~$\Mlayer$ --- census устойчивых составных конфигураций с нетривиальным
топологическим/цветовым бюджетом.

\section{Конфайнмент без нового поля}
\label{sec:floor2-confinement}

«Вытащить кварк'' на макроязыке --- разорвать составной топологический узел.
На~$\Mlayer$ это не извлечение одного планкона из ячейки, а
\emph{пара} (или больше) планконов с разнесёнными фокусами:
энергия растёт с расстоянием, потому что ломается holonomy и занятость на звезде~$\varepsilon$.

Цветовой заряд $SU(3)$ \emph{не вводится} отдельной калибровкой сверху.
Ожидание: кратность устойчивых фокусов в барионном/мезонном узле
следует из топологии на FCC/$\varepsilon$ + Паули +~$K_P$;
естественный порядок величины --- $\sim d=3$ (три пространственных измерения emergent 3D,
\cref{ch:carrier}).

Имя $N_c=3$ в Стандартной модели --- язык~$T$; на~$\Mlayer$ первичны census и winding,
не таблица PDG.

\section{Квантовые числа}
\label{sec:floor2-quantum-numbers}

На этаже~2 впервые появляются (или становятся наблюдаемыми) составные квантовые числа,
которые на этаже~0 были только зародышами:

\begin{table}[htbp]
  \centering
  \caption{Откуда берутся квантовые числа (схема).}
  \label{tab:floor2-qnumbers}
  \begin{tabular}{@{}lll@{}}
    \toprule
    число & этаж~0 & этаж~2 \\
    \midrule
    электрический $Q$ & $n\in\mathbb{Z}$ на контуре & сохраняется в составном узле \\
    спин & $SU(2)$ на ядре, $2\pi\to -1$ & сумма / связка фокусов \\
    цвет & --- & фокусы на $\varepsilon$-звезде, $\sim d$ \\
    поколение & --- & census стеков / иерархия $\alpha$ (открыто) \\
    барионное $B$ & $b$ на ячейке & составной бюджет узла \\
    \bottomrule
  \end{tabular}
\end{table}

Полный вывод таблицы --- не закрыт: это программа этажа~2,
а не сводка готовой теории.

\section{Связь с сильным взаимодействием}
\label{sec:floor2-strong}

На этаже~2 сильное взаимодействие --- не отдельный Лагранжиан,
а \emph{короткий} канал packing/confining той же фазовой башни,
что даёт электрослабое сопряжение на больших $N$
(\cref{ch:alpha,ch:si-sm}).

Шкала семени иерархии $\alpha$ совпадает с электрослабым $v$:
$N_{\mathrm{hier}}=8$ ступеней ослабления на $\lfloor B_{hV}\rfloor-1$.
Сильный канал --- нижняя ступень той же башни, не новый угол~$\pi$.

\section{Открыто на этаже~2}
\label{sec:floor2-open}

\begin{itemize}[noitemsep]
\item полный census составных узлов (барионы, мезоны, тяжёлые резонансы);
\item явный sim конфайнмента и спектра мод на $\varepsilon$-решётке;
\item поколения и CKM-углы как следствие census, не как вход;
\item стыковка с $k=7\ldots 9$ IR-лестницы ($r_e$, $\bar\lambda_C$, $a_0$).
\end{itemize}

После этажа~2 лестница выходит на инфракрасную ось
$N_k=N_c/\alpha^{k}$ (\cref{eq:ir-hop-ladder}): там начинается атомная и макро-физика на~$T$.

\paragraph{Что читать дальше.}
Константы связи и массы --- \cref{ch:alpha,ch:si-sm};
гравитация и electroweak как readout той же среды --- в главах о силах модели
(см. оглавление).
